\section{Walk the Keys, Not the Rows}
\label{sec:ch11-walk-the-keys-not-the-rows}

\index{DataLoader!the corrected loop|(}%
The fix is the loop, rewritten to count the other collection. Here is
\code{SpeakerDataLoader.cs} again, complete, with what moved marked:

\begin{minted}[highlightlines={13-17,40,42-49}]{csharp}
using GreenDonut;
using Microsoft.EntityFrameworkCore;

namespace Speakers;

/// <summary>
/// What the source-generated loader is underneath: a class deriving
/// from DataLoaderBase, given the batch of keys and a span to put the
/// answers in. Written out by hand because a store that is not a
/// DbContext - a remote service, a cache, a stored procedure - hands
/// back a list rather than a dictionary, and the generator only
/// accepts a dictionary.
///
/// The whole of the difference between this and the generated one is
/// which collection the loop walks. FetchAsync's own doc comment says
/// it: for every provided key must be a result, in the exact same
/// order the keys were provided.
/// </summary>
public interface ISpeakerByIdDataLoader : IDataLoader<int, Speaker>;

public sealed class SpeakerByIdDataLoader(
    IServiceProvider services,
    IBatchScheduler batchScheduler,
    DataLoaderOptions options)
    : DataLoaderBase<int, Speaker>(batchScheduler, options),
        ISpeakerByIdDataLoader
{
    protected override async ValueTask FetchAsync(
        IReadOnlyList<int> ids,
        Memory<Result<Speaker?>> results,
        DataLoaderFetchContext<Speaker> context,
        CancellationToken ct)
    {
        await using var scope = services.CreateAsyncScope();
        var db = scope.ServiceProvider
            .GetRequiredService<SpeakerContext>();

        var rows = await db.Speakers
            .Where(s => ids.Contains(s.Id))
            .ToDictionaryAsync(s => s.Id, ct);

        // Walk the keys, not the rows. The span is positional and the
        // position that matters is the key's, so every slot is written
        // exactly once and a key with no row is written as a null
        // rather than left for the next row to land in.
        for (var i = 0; i < ids.Count; i++)
        {
            results.Span[i] = Result<Speaker?>.Resolve(
                rows.GetValueOrDefault(ids[i]));
        }
    }
}
\end{minted}

\index{DataLoader!the loop bound is the tell}%
\code{ToDictionaryAsync} instead of \code{ToListAsync}, and the loop counts to
\code{ids.Count} rather than \code{rows.Count}. If you want one thing to look
for when reading somebody else's loader, that bound is it. A loop over the
results is a loop that has decided the store's order is the caller's order, and
it is wrong even when it works.

\index{DataLoader!GetValueOrDefault does the nulls}%
\code{GetValueOrDefault} is doing the other clause. A key with no row yields a
null, which lands in that key's own slot, which is the null the specification
asks for and the one chapter~\ref{ch:router-n-plus-1} showed a missing session
producing.

\index{DataLoader!indistinguishable from the generated one}%
With that loop, every request in this chapter answers the way the
source-generated loader answers it. The two sessions come back carrying their
own speakers. The out-of-key-order list of three comes back in the order it was
asked for. The middle key that matches no row is a null in its own position
rather than an error. Every count in
chapters~\ref{ch:second-third-subgraph} and~\ref{ch:router-n-plus-1} is the
number those chapters printed, and the statement is the same batched statement
with the same three keys in it.

\subsection{The Generator Was Already Doing This}

\index{DataLoader!what the source generator emits}%
The reason chapters~\ref{ch:data-without-n-plus-1}
to~\ref{ch:router-n-plus-1} never met this problem is that
the \code{[DataLoader]} attribute writes that loop for you, and writes it the
same way every time. Given a method returning a dictionary, the generator emits a
\code{CopyResults} that walks the keys, calls \code{TryGetValue} on each one,
and puts a default in the slot when the lookup misses. It cannot do anything
else, because a dictionary has no order to be tempted by.

\index{DataLoader!a dictionary cannot be misordered}%
Which turns out to be a better argument for the source generator than
convenience is. A loader that hands back a dictionary is a loader that cannot
make this mistake, and the shape of the return type is what rules it out. I
would not have written that sentence before this chapter; I used
the attribute because it saved me writing the file.

\index{DataLoader!a list return is not a batch loader}%
The obvious middle road is closed, and it closes badly. Give a
\code{[DataLoader]} method the return type \code{Task<IReadOnlyList<Speaker>>},
expecting the generator to map that list onto the keys, and it does not reject
the method. It reinterprets it: the key becomes the whole
\code{IReadOnlyList<int>}, the value becomes the whole list of speakers, and
what comes out is a cache loader that fetches one at a time. The call site then
fails to compile, twice, with a message about converting \code{object} to
\code{Speaker} that mentions none of this. Caught, which is what matters, and
caught by the type checker rather than by anything that knows what you meant.

\subsection{The Assertion Worth Writing}

\index{testing!ask the entity route for its own key}%
There is one test for this and it does not depend on knowing anything about
your data. Ask a subgraph's entity route for the key alongside the fields, in
an order that is not key order, and check that each entry carries the key its
own representation named:

\begin{minted}[fontsize=\footnotesize,breakanywhere]{json}
{"r":[{"__typename":"Speaker","id":"U3BlYWtlcjoz"},
      {"__typename":"Speaker","id":"U3BlYWtlcjox"},
      {"__typename":"Speaker","id":"U3BlYWtlcjoy"}]}
\end{minted}

\begin{minted}[fontsize=\footnotesize,breakanywhere]{json}
{"data":{"_entities":[
  {"id":"U3BlYWtlcjoz","name":"Chidi Okafor"},
  {"id":"U3BlYWtlcjox","name":"Ada Fischer"},
  {"id":"U3BlYWtlcjoy","name":"Bruno Kaminski"}]}}
\end{minted}

\index{testing!the answer carries its own evidence}%
The evidence is inside the answer. Every entry names a key, so a machine can
compare it with the key that entry's representation carried, and it can do that
without knowing which speaker is which or what order the store prefers. Against
the broken loader the same request comes back with the three keys ascending,
so the first entry says \code{U3BlYWtlcjox} where the representation asked for
\code{U3BlYWtlcjoz}. Every name is still attached to its own key; what has gone
wrong is which question each pair is answering, and asking for the key is what
makes that visible.

\index{testing!the router never asks for the key}%
That works because you asked for the key, which is exactly what the router does
not do. There is no way to make the router perform this check, and I would not
want the column on every entity fetch in production even if there were. It is a
test, run against the subgraph's own port, and it belongs beside the tests
chapter~\ref{ch:testing} builds.

\index{testing!an ascending fixture proves nothing}%
Write the fixture in an order the store will not reproduce, and it is worth
being deliberate about that rather than lucky. A fixture that happens to list
its keys ascending passes against both loaders in this chapter and tells you
nothing at all, which is the same trap the whole graph fell into in
section~\ref{sec:ch11-why-this-graph-got-away-with-it}, one layer down.

\subsection{And When You Cannot Write One}

\index{DataLoader!the remote store case}%
The case that got me into this file is worth naming, because it is the one that
takes people out of the source generator in the first place. A loader that
calls another service asks for a list of ids over HTTP and is handed a JSON
array. There is no dictionary, there is nobody's primary key to sort by, and
the service on the other end has made no promise about sequence. Build the
dictionary yourself, off whatever the payload calls its identifier, before you
go anywhere near the span. The loop stays the same loop.

\index{DataLoader!the store that answers a different length}%
And check the length rather than trusting it. A remote service that drops the
ids it could not find gives you a shorter array than you asked for, which is
the first clause of the contract breaking, and the router will refuse the
answer with a message about the returned entities count not matching. That is
the good outcome. The bad one is a service that also happens to return them
sorted, where the count is right and nothing is.
\index{DataLoader!the corrected loop|)}%
